\subsection{Milestone 2 -- Identificazione di modelli accurati}

I risultati completi delle 60 combinazioni sono nelle Tabelle~\ref{tab:r2-rf-full},
\ref{tab:r2-nb-full}, \ref{tab:r2-ibk-full}.

\subsubsection{Kappa e dataset sbilanciato}

Con il 9\% di istanze buggy (v.\ Sez.~2.2.3), l'accuracy è fuorviante e Kappa
è la metrica discriminante.

RF assegna a ogni classe un punteggio da 0 a 1 che rappresenta la probabilità stimata
di essere buggy. AUC $= 0{,}94$ significa che, presa una classe buggy e una clean a
caso, nel 94\% dei casi la buggy riceve un punteggio più alto. Kappa $= 0{,}61$ misura
quanto le predizioni binarie finali superano il caso: su un dataset al 9\% di buggy,
predire sempre ``clean'' dà Kappa $\approx 0$; ottenere $0{,}61$ significa che il modello
classifica correttamente molto più di quanto ci si aspetterebbe per caso.

\subsubsection{Effetto del bilanciamento}

Undersampling aumenta il Recall di RF ($0{,}52 \to 0{,}86$) ma fa crollare Kappa
($0{,}58 \to 0{,}37$) e Precision ($0{,}73 \to 0{,}31$): il training 50/50 spinge a predire
``buggy'' più spesso del necessario, gonfiando i falsi positivi. Oversampling e SMOTE ottengono un
bilancio migliore; SMOTE è superiore perché genera istanze sintetiche per interpolazione
invece di duplicare, aumentando la variabilità della classe minoritaria in modo più granulare.

\subsubsection{NaiveBayes insensibile al bilanciamento}

NB guadagna solo $+3$\, di Recall con Undersampling, contro $+34$\, di RF; Kappa rimane
$0{,}26$--$0{,}27$ indipendentemente dalla strategia di bilanciamento. Il bilanciamento non risolve
il problema strutturale di NB: LOC\_Touched, Churn e LOC\_Added sono fortemente correlate
tra loro, e NB le tratta come indipendenti.

\subsubsection{RandomForest domina AUC e NPofB20}

RF produce probabilità come media su 100 alberi, ottenendo punteggi continui
e ben calibrati per il ranking delle istanze. IBk con $k=1$ non produce
punteggi graduati ma assegna direttamente la classe dell'istanza più vicina,
rendendo il ranking grossolano; NB, non soddisfacendo l'assunzione di
indipendenza a causa della correlazione tra feature, produce stime di
probabilità distorte. Il risultato migliore è: NPofB20 di RF $= 0{,}88$.

\subsubsection{Feature selection: filtri ininfluenti, wrapper controproducenti}

I metodi filter (InfoGain, Spearman) risultano ininfluenti su questo dataset: tutte le 19
feature hanno InfoGain $> 0$, quindi \texttt{Ranker(threshold=0.0)} non ne elimina nessuna.
La differenza di $0{,}01$ su Recall tra No FS+SMOTE ($0{,}66$) e Spearman/InfoGain+SMOTE
($0{,}67$) è rumore della CV: le tre configurazioni usano esattamente le stesse feature.
I wrapper (ForwardSearch, BackwardSearch) selezionano il sottoinsieme ottimizzando per
NaiveBayes, non per RF: il risultato peggiora RF rispetto a No FS e produce esiti misti
per IBk (ForwardSearch+SMOTE massimizza AUC a $0{,}80$ ma Precision scende a $0{,}38$).

\subsubsection{Scelta del modello migliore}

\textbf{RF + No FS + SMOTE} emerge come configurazione migliore per uso generale, con il
miglior bilanciamento complessivo tra le cinque metriche:

\begin{itemize}[noitemsep]
  \item \textbf{Kappa $0{,}61$}: NaiveBayes non supera $0{,}29$ per la violazione
        dell'indipendenza condizionale; IBk si ferma a $0{,}47$ per la dipendenza da un
        singolo vicino con 19 feature.
  \item \textbf{Precision $0{,}63$}: RF + Undersampling scende a $0{,}31$, rendendo
        il modello inutilizzabile in pratica.
  \item \textbf{Recall $0{,}66$}: RF senza balancing si ferma a $0{,}52$ perché predire
        quasi sempre ``clean'' minimizza già l'errore globale.
  \item \textbf{AUC $0{,}94$} e \textbf{NPofB20 $0{,}88$}: nettamente superiori a NB
        e IBk per la qualità dei punteggi interni di RF.
\end{itemize}

Kappa e NPofB20 hanno pesato di più nella scelta finale. Kappa è la metrica più onesta su
un dataset al 9\% di buggy: corregge lo sbilanciamento e rende il confronto con la baseline
``predici sempre clean'' significativo. NPofB20 è direttamente rilevante nella pratica: misura
quanti bug si trovano ispezionando solo le classi più rischiose, che è esattamente come il
risultato verrebbe usato in un contesto reale. SMOTE è preferito a Oversampling per un margine
stretto su entrambe le metriche (Kappa $0{,}61$ vs $0{,}59$, NPofB20 $0{,}88$ vs $0{,}87$).
No FS è la scelta più semplice: nessun passo aggiuntivo, nessun peggioramento.
